\documentclass[11pt]{article}
\usepackage[margin=1in]{geometry}
\usepackage{amsmath,amssymb,mathtools}
\usepackage{fontspec,unicode-math}
\setmainfont{Latin Modern Roman}
\setsansfont{Latin Modern Sans}
\setmathfont{Latin Modern Math}
\usepackage{xurl,hyperref}
\hypersetup{colorlinks=true,linkcolor=blue,urlcolor=blue}
\newcommand{\sref}[2]{\hyperlink{source:#1:#2}{[S#2]}}
\newcommand{\cataloguescope}[1]{\paragraph{Mathematical question.}#1}
\begin{document}
\section{Newton dynamics for entire functions}\label{statement}

\subsection{Definitions and mathematical statement}

\paragraph{Shared notation.}$\mathbb N=\{1,2,\ldots\}$, and $\mathbb Z,\mathbb Q,\mathbb R,\mathbb C$ denote the integers, rationals, reals and complex numbers. Any other conventions are specified locally.


For a nonconstant entire function $f$, Newton iteration uses the meromorphic map
\[N_f(z)=z-\frac{f(z)}{f'(z)}.\]
At poles the map is considered in the Riemann sphere; its dynamics concerns the forward iterates where these are defined. A transcendental entire function is an entire function that is not a polynomial.
\paragraph{Statement recorded in the supplied source.}
Describe the dynamics of Newton iteration for general classes of transcendental entire functions. The target is a general dynamical description, allowing the choice and analysis of natural classes of functions. \sref{statement}{2}
\subsection{Short English statement}
Develop a dynamical description of Newton’s method for general classes of transcendental entire functions.
\subsection{Sources}
\begin{description}
\item[\hypertarget{source:statement:1}{S1}] ULAM.AI and the UnsolvedMath Contributors, UnsolvedMath: A Curated Collection of Open Mathematics Problems, catalogue entry. CC BY 4.0. \url{https://huggingface.co/datasets/ulamai/UnsolvedMath}.
\item[\hypertarget{source:statement:2}{S2}] Ben Bielefeld, Mikhail Lyubich and contributors, \emph{Problems in Holomorphic Dynamics}, arXiv:math/9205209v1 (1992), Robert Devaney, ``Open Questions in Non-Rational Complex Dynamics'', \S5, eighth Problem, printed p. 36. \url{https://arxiv.org/abs/math/9205209v1}.
\end{description}
\label{end:statement}
\end{document}
